\section{Results}\label{sec:results}

We report results from the released corpus of $4{,}860$ TDMA-CBBA trials ($324$ configurations, $15$ independent seeds each; Sec.~\ref{sec:test-matrix}) together with the four-baseline ablation sweep ($900$ runs) and the communication--motion isolation ablation of Sec.~\ref{sec:setup-ablation}.
All figures are generated directly from raw CSV by the open-source analysis pipeline; no manual curation is involved.

\subsection{Headline Statistics}
Table~\ref{tab:headline} summarizes the four metrics across the full corpus.
Every metric rejects normality (Shapiro--Wilk, all $p < 10^{-4}$, $n = 4860$), justifying the non-parametric and survival treatments of Sec.~\ref{sec:setup-stats}.

\begin{table}[htbp]
\centering
\caption{Headline metrics over 4{,}860 TDMA-CBBA trials (mean $\pm$ SEM, bootstrap 95\% CI). MTC is right-censored at 50\,s.}
\label{tab:headline}
\begin{tabular}{@{}lccc@{}}
\toprule
Metric & Mean $\pm$ SEM & 95\% CI \\
\midrule
MTC (s) & $21.57 \pm 0.315$ & $[20.95,\ 22.21]$ \\
SCI     & $0.577 \pm 0.002$ & $[0.574,\ 0.581]$ \\
CE      & $0.644 \pm 0.004$ & $[0.636,\ 0.652]$ \\
PDR     & $0.827 \pm 0.002$ & $[0.823,\ 0.832]$ \\
\midrule
Convergence rate & $65.1\%$ & --- \\
\bottomrule
\end{tabular}
\end{table}

The aggregate PDR of $0.938$ is consistent with the two-stage loss model of Sec.~\ref{sec:method-loss} averaged over the loss sweep, providing an end-to-end sanity check that the radio simulator behaves as designed (see Sec.~\ref{sec:results-pdr}).

\subsection{Convergence Behavior}
Fig.~\ref{fig:mtc_vs_swarm} shows MTC against swarm size.
Consensus time grows approximately linearly with $n$ at low loss, consistent with the scheduling-delay term $N\tau K$ in Theorem~\ref{thm:noiseless}: each additional drone both lengthens the frame and adds one ledger to synchronize.
At larger swarms and higher loss rates the curve separates into a converging branch and a saturated branch pinned at the 50\,s horizon---the two regimes delimited by Theorems~\ref{thm:noiseless} and \ref{thm:impossibility}, with Theorem~\ref{thm:probabilistic} bounding the expiry-free probability of the deterministic precondition.

\begin{figure}[htbp]
\centering
\includegraphics[width=\columnwidth]{mtc_vs_swarm_size.png}
\caption{Mean time to consensus vs.\ swarm size (error bars: std.\ across trials). Beyond the staleness threshold $N\tau > T_{\mathrm{val}} = 2$\,s of Thm.~\ref{thm:noiseless}(ii), MTC saturates at the mission horizon, marking the synchronization--staleness phase transition.}
\label{fig:mtc_vs_swarm}
\end{figure}

The convergence CDF (Fig.~\ref{fig:convergence_curve}) makes this bimodality explicit: the empirical distribution exhibits a steep rise at short times (sub-threshold trials) followed by a long flat plateau (censored trials above threshold), rather than the unimodal tail a purely stochastic slowdown would produce.

\begin{figure}[htbp]
\centering
\includegraphics[width=\columnwidth]{convergence_curve.png}
\caption{Empirical CDF of time to consensus. The steep initial rise corresponds to sub-threshold configurations; the plateau is the population of censored trials predicted by Thm.~\ref{thm:noiseless}(ii).}
\label{fig:convergence_curve}
\end{figure}

\subsection{Packet Delivery Ratio}\label{sec:results-pdr}
Fig.~\ref{fig:pdr_vs_pl} shows PDR degrading gracefully with nominal transmit-leg loss.
The observed per-link delivery ratios track the joint end-to-end
success probability $(1-p_{\mathrm{tx}})(1-p_{\mathrm{rx}})$ predicted
by the two-stage model of Sec.~\ref{sec:method-loss}: with the revised
offered-link accounting (which counts every scheduled sender--receiver
opportunity including transmit-side failures) the measured ratios
($0.927$, $0.787$, $0.594$ at $5\%$, $15\%$, $30\%$ loss) agree with
the closed-form estimand to within $10^{-3}$ (see
\texttt{pdr\_model\_check.txt}).

\begin{figure}[htbp]
\centering
\includegraphics[width=\columnwidth]{pdr_vs_packet_loss.png}
\caption{Packet delivery ratio vs.\ nominal transmit-leg packet loss (error bars: 95\% CI across trials).}
\label{fig:pdr_vs_pl}
\end{figure}

\subsection{Distributional Structure}
The heatmap (Fig.~\ref{fig:heatmap}) locates the phase transition of Theorem~\ref{thm:noiseless}(ii) in parameter space: mean MTC rises smoothly in the low-$n$, low-loss corner, then jumps discontinuously toward the horizon cap as $N\tau$ approaches and exceeds $T_{\mathrm{val}} = 2$\,s---the predicted boundary running diagonally through the $(n, p)$ plane.

\begin{figure}[htbp]
\centering
\includegraphics[width=\columnwidth]{heatmap.png}
\caption{Mean MTC over swarm size $\times$ packet loss, annotated with convergence rate. The jump toward the cap traces the synchronization--staleness phase transition of Thm.~\ref{thm:noiseless}(ii).}
\label{fig:heatmap}
\end{figure}

\subsection{Censoring-Aware Survival Analysis}\label{sec:results-survival}
Because consensus is not observed in $34.9\%$ of trials, we treat time-to-consensus as right-censored survival data rather than a numerical scalar.
Fig.~\ref{fig:survival} shows Kaplan--Meier survivor curves per baseline: ideal and loss-only open-channel swarms converge almost immediately (median $0.2$--$0.4$\,s), TDMA-only converges with a moderate tail (median $3.3$\,s), while TDMA-with-loss retains a $64.9\%$ convergence fraction (median $6.4$\,s), its survivor curve never reaching the baseline---the signature of the structural phase transition rather than a slow-but-certain process. Pairwise log-rank tests and a TDMA-only Cox proportional-hazards fit (staleness ratio and packet loss as covariates, robust standard errors) supplement the curves (Sec.~\ref{sec:setup-stats}); full outputs are in \texttt{survival\_analysis.txt}.

\begin{figure}[htbp]
\centering
\includegraphics[width=\columnwidth]{survival_curves.png}
\caption{Kaplan--Meier survival curves for time to consensus, per baseline. Right-censoring at 50\,s is respected (drops = failures to converge).}
\label{fig:survival}
\end{figure}

\subsection{Factorial Regression and Ablation Evidence}\label{sec:results-regression}
Within TDMA conditions we use regression as a descriptive mechanism analysis rather than a causal estimator (Eq.~\eqref{eq:regression}, Sec.~\ref{sec:setup-stats}): the staleness ratio $\rho = N\tau/T_{\mathrm{val}}$ and the packet-loss rate are treated as controlled experimental factors, and HC3 robust standard errors are reported. The estimated interaction $\beta_3$ between $\rho$ and $p_{\mathrm{tx}}$ is consistent with the joint scheduling-and-loss mechanism of Thm.~\ref{thm:noiseless}(ii); full coefficient tables are in \texttt{regression\_analysis.txt}.

Table~\ref{tab:ablation} summarizes the four-baseline sweep, which refines the causal decomposition of Sec.~\ref{sec:baselines}: moving from ideal (MTC $0.34$\,s) to TDMA-only raises MTC to
$3.46$\,s (converging $100\%$); adding loss pushes TDMA to $22.4$\,s
/ $61.4\%$. Crucially, open-channel loss-only converges $100\%$ at every loss level---fresh every-frame broadcasting defeats information expiry, isolating belief staleness as the mechanism that breaks TDMA operation.

\begin{table}[htbp]
\centering
\caption{Four-baseline ablation ($N \in \{10,20,30,50\}$). MTC is right-censored at 50\,s.}
\label{tab:ablation}
\begin{tabular}{@{}lccc@{}}
\toprule
Baseline & MTC (s) & Convergence & PDR \\
\midrule
Ideal CBBA            & $0.34$ & $100\%$   & $1.000$ \\
Loss-only (open ch.)  & $0.47$ & $100\%$   & $0.829$ \\
TDMA-only             & $3.46$ & $100\%$   & $1.000$ \\
TDMA-CBBA (proposed)  & $22.4$ & $61.4\%$  & $0.828$ \\
\bottomrule
\end{tabular}
\end{table}

Finally, we run the communication--motion isolation ablation of Sec.~\ref{sec:setup-ablation} (motion disabled, all four channel models) to separate auction-communication effects from stale flocking/collision state; the per-configuration results are released with the raw data and indicate that the MTC ordering across baselines is preserved when motion is disabled, confirming the communication layer as the dominant driver of the reported differences.
% placeholder